\chapter*{Test: Min System Call Latency}{\textbf{Section Authors: Chenyi Wang}}
 \newline 
 \newline The goal of the Chapter is to compare AOS with xv6 open source system. We measured how long the cheapest possible system call takes on each one.

\addcontentsline{toc}{chapter}{Test: System Call Latency}

\section{Set up}
We use the \co{https://github.com/k-mrm/xv6-aarch64?tab=readme-ov-file} as the baseline of performance. And we added the \co{syscall1(SYSCALL_NOP)} in the \co{hello.c} and repeated for 100000 times which is to show the cost of min syscall. And on xv6 the syscall in the loop is \co{getpid()}.

\section{Results}

The counter on this host runs at 62.5 MHz, so one tick is sixteen nanoseconds. Table \ref{tab:syscall-latency} shows the per-call cost on each system and the implied wall clock time.

\begin{table}[h]
\centering
\begin{tabular}{lrrr}
\hline
System & Syscall used & Ticks per call & Time per call \\
\hline
xv6 & \texttt{getpid} & 115 & 1.84 \(\mu\)s \\
AOS & \texttt{SYSCALL\_NOP} & 408 & 6.53 \(\mu\)s \\
\hline
\end{tabular}
\caption{Per-call cost of the cheapest system call on each system,
measured on QEMU with one CPU.}
\label{tab:syscall-latency}
\end{table}

\subsection{Analysis}

A syscall on AOS is about 3.5x times more expensive than xv6. This matches what we would expect from the
two designs. xv6 follows the usual monolithic Unix pattern: the trap goes straight to a syscall number, looks the number up in a small table, and runs the handler. 
AOS, being a multikernel derived from Barrelfish, shares a single trap entry across several very different request types. Even when the actual handler does nothing, the kernel still has
to decode a packed \co{union syscall_info} word.